\chapter{Implementation}
\label{ch:implementation}

This chapter walks through the code that turns the design in \Cref{ch:design} into a running system. It is organised top-down: server bootstrap, transport and middleware, the auth subsystem, the exam aggregate, the auto-grader, the manual-grading pipeline, the reminder scheduler, the React SPA, and a short pass through the Postgres-side conventions. Where a code listing is shown inline, its enclosing function and the line range from the source file are cited; long listings are pushed to \Cref{appendix:source-listings}.

\section{Backend bootstrap}

The Go process is configured by environment variables (\texttt{config.Load}) and started by \texttt{main.go:27--56}. The order is fixed: load config; connect Postgres (which runs migrations); initialise the auth, email, and Judge0 packages; start the reminder goroutine; build the Gin router; register routes; serve. \Cref{fig:F5.1} shows the call graph.

\begin{figure}[H]
  \centering
  \includegraphics[width=\linewidth]{F5.1_bootstrap-callgraph.pdf}
  \caption{Bootstrap call graph from \texttt{main.go:27--56}.}
  \label{fig:F5.1}
\end{figure}

\texttt{database.Connect} (\texttt{internal/database/database.go}) opens the Postgres connection, runs \texttt{RunMigrations}, then \texttt{AutoMigrate}, then \texttt{RunPostMigrations}. The split is documented in \Cref{ch:design} and codified in \texttt{ORM\_RULES.md}: idempotent SQL backfills run \emph{before} \texttt{AutoMigrate} so they can refer to columns and tables that already exist, and \texttt{RunPostMigrations} runs \emph{after} \texttt{AutoMigrate} so it can reference tables that GORM just created (e.g.\ adding a \texttt{folders} foreign key on \texttt{problems.folder\_id}).

The router is split into \texttt{public} and \texttt{protected} groups (\texttt{main.go:38--40}). \texttt{public} is open; \texttt{protected} is wrapped in \texttt{middleware.Auth(cfg.JWTSecret)}, which means every protected handler can rely on \texttt{c.MustGet("user\_id")} and \texttt{c.MustGet("role")} returning a sensible value.

\section{Transport: middleware chain}

Two pieces of middleware are interesting: the JWT authenticator and the role guard. \Cref{fig:F5.2} draws the chain.

\begin{figure}[H]
  \centering
  \includegraphics[width=0.9\linewidth]{F5.2_middleware-chain.pdf}
  \caption{Middleware chain. CORS handles preflight; \texttt{Auth} validates and unpacks the token; \texttt{RequireRole} guards role-restricted route groups; per-resource ownership re-checks happen inside handlers.}
  \label{fig:F5.2}
\end{figure}

\Cref{lst:jwt-parse} is the heart of \texttt{middleware.Auth} (\texttt{internal/middleware/auth.go:20--56}). The two details to flag are the signing-method check (anything other than HMAC is rejected) and the comma-ok type assertions on the claims. Earlier revisions of this file used unchecked assertions (\texttt{userIDFloat := claims["user\_id"].(float64)}); a malformed token then crashed the process via panic. The fix landed in commit \texttt{2eaac88} along with the rest of the security review.

\begin{figure}[H]
\begin{lstlisting}[language=Go,caption={JWT parse with safe claim extraction (\texttt{internal/middleware/auth.go:20--56}).},label={lst:jwt-parse}]
token, err := jwt.Parse(tokenStr, func(t *jwt.Token) (any, error) {
    if _, ok := t.Method.(*jwt.SigningMethodHMAC); !ok {
        return nil, jwt.ErrSignatureInvalid
    }
    return []byte(secret), nil
})
if err != nil || !token.Valid {
    c.AbortWithStatusJSON(http.StatusUnauthorized, gin.H{"error": "invalid token"})
    return
}
claims, ok := token.Claims.(jwt.MapClaims)
if !ok { /* 401 */ return }

userIDFloat, ok := claims["user_id"].(float64)
if !ok { /* 401 */ return }
email, ok := claims["email"].(string)
if !ok { /* 401 */ return }
role, ok := claims["role"].(string)
if !ok { /* 401 */ return }

c.Set("user_id", uint(userIDFloat))
c.Set("email", email)
c.Set("role", role)
c.Next()
\end{lstlisting}
\end{figure}

\texttt{middleware.RequireRole} (\texttt{internal/middleware/auth.go:61--81}) is a one-page allow-list helper: given one or more role strings, it admits the request if \texttt{c.Get("role")} matches any of them, and otherwise responds 403. Role checks happen \emph{after} authentication, so the helper assumes a role exists in the context and treats absence as 401.

\section{Auth subsystem}

The auth handlers live in \texttt{internal/auth/handler.go}. They cover four flows: signup/login with bcrypt, password reset by email, Google OAuth via OpenID Connect, and full-account deletion.

\paragraph{Signup and login.} \texttt{auth.Signup} (\texttt{internal/auth/handler.go:43--97}) validates email format with a small regex, requires at least 6-character passwords (production deployments will want a stronger minimum; this is the lower bound), hashes with bcrypt at the default cost~\cite{provos1999bcrypt}, and creates the matching \texttt{user\_profiles} row with notification toggles defaulted to on. \texttt{auth.Login} (\texttt{internal/auth/handler.go:99--128}) is symmetric: a missing user, an empty hash, or a bcrypt mismatch all produce the same generic 401 message so the endpoint cannot be used as an account enumeration oracle.

\paragraph{Password reset.} The reset flow is documented in \Cref{ch:security} (\Cref{fig:F6.5}). The interesting choice is that \texttt{ForgotPassword} (\texttt{internal/auth/handler.go:289--340}) always returns the same generic success message regardless of whether the email exists, again to avoid an enumeration oracle. The token is generated as 32 random bytes (\texttt{crypto/rand.Read}), hex-encoded for the link, and stored as its SHA-256 hash in \texttt{password\_reset\_tokens.token\_hash}. Storing the hash means a Postgres dump cannot be replayed against the live system. Prior unused tokens for the same user are invalidated by setting their \texttt{used\_at} to "now" before inserting the new one --- this is what makes the token strictly single-shot even if a user clicks "forgot password" twice.

\paragraph{Google OAuth.} \texttt{auth.GoogleAuth} (\texttt{internal/auth/handler.go:389--473}) accepts an ID token from the SPA, validates it through Google's official Go SDK (\texttt{idtoken.Validate} from \texttt{google.golang.org/api/idtoken}), and dispatches across three cases: match by \texttt{google\_id}, link to an existing email, or create a new account once a role has been provided. The flow guards on \texttt{email\_verified}: an unverified Google email is rejected with 401. Newly-created Google accounts have an empty \texttt{password\_hash}; \texttt{auth.Login} explicitly refuses to authenticate empty hashes, so a Google account cannot be hijacked by guessing the password.

\paragraph{Account deletion.} \texttt{auth.DeleteAccount} (\texttt{internal/auth/handler.go:130--273}) does the whole job in a single transaction. The student path deletes the user's class memberships, the test results that reference their submissions, the submissions themselves, and their exam attempts. The teacher path is heavier: it walks the teacher's classes and exams, and recursively deletes the test results, submissions, test cases, problems, exam attempts, exam-class links, exams, announcements, and classes. The transaction rollback path is exhaustive --- every error returns and rolls back --- so a partial deletion cannot leave dangling references. This is the path that motivated dropping the legacy \texttt{messages}, \texttt{teams}, and \texttt{team\_members} tables: their FK constraints on \texttt{users(id)} blocked account deletion until they were removed.

\section{Exam aggregate}

\texttt{internal/exam/handler.go} is the largest file in the backend (751 lines). It owns the teacher CRUD surface for exams, the assignment endpoint, the close/reopen endpoints, and the three student endpoints \texttt{StartAttempt}, \texttt{SubmitAttempt}, and \texttt{AutosaveAttempt}. \texttt{internal/exam/reset.go} adds an extra teacher action: destructive reset (clearing student attempts) when an exam's content has changed materially after publishing.

\paragraph{Create and update.} \texttt{exam.CreateExam} (\texttt{internal/exam/handler.go:32--82}) defaults a new exam to \texttt{is\_draft = true}, applies sensible duration and passing-score defaults, and parses the optional \texttt{start\_time} / \texttt{end\_time} via a small list of accepted layouts (\texttt{exam.parseTime}, \texttt{internal/exam/handler.go:737--751}). \texttt{exam.UpdateExam} (\texttt{internal/exam/handler.go:179--237}) is a partial update, but with one absolute rule: publishing requires a \texttt{start\_time}.

\begin{figure}[H]
\begin{lstlisting}[language=Go,caption={Publish-with-schedule check (\texttt{internal/exam/handler.go:222--233}).},label={lst:publish-check}]
publishing := req.IsDraft != nil && !*req.IsDraft
willHaveStart := exam.StartTime != nil
if t, ok := updates["start_time"].(time.Time); ok {
    willHaveStart = !t.IsZero()
}
if publishing && !willHaveStart {
    c.JSON(http.StatusBadRequest, gin.H{"error": "set a start time before publishing"})
    return
}
\end{lstlisting}
\end{figure}

\paragraph{Assignment.} \texttt{exam.AssignExam} (\texttt{internal/exam/handler.go:352--388}) replaces the set of \texttt{exam\_classes} rows for an exam. It first verifies that every requested class belongs to the calling teacher (\texttt{COUNT(*)} matches \texttt{len(req.ClassIDs)}), then deletes the old links and inserts the new ones. The "delete + insert" path is intentionally not transactional: the worst case is a brief window during which an exam is unassigned, and assignment changes are rare enough that an attempt to take advantage of it is implausible. If that ever changes, wrapping the two operations in \texttt{database.DB.Transaction} is a one-line edit.

\paragraph{Starting an attempt.} \texttt{exam.StartAttempt} (\texttt{internal/exam/handler.go:468--530}) is idempotent. It loads the exam, verifies that the caller is enrolled in at least one class to which the exam is assigned, refuses to start a draft, and refuses outside the time window. If an attempt for \texttt{(user\_id, exam\_id)} already exists, it is returned; otherwise a new one is created. The unique index on those two columns means two concurrent calls cannot create duplicate attempts.

\paragraph{Submitting and finalising.} \texttt{exam.SubmitAttempt} (\texttt{internal/exam/handler.go:547--643}) creates one \texttt{Submission} row per answer, each linked to the attempt via \texttt{exam\_attempt\_id}, and dispatches to one of three graders by \texttt{Problem.Type}. Coding submissions go to \texttt{judge0.Grade}, MCQs to \texttt{judge0.GradeMCQ}, written answers to \texttt{judge0.GradeWritten}. The handler then marks the attempt \texttt{submitted}, clears \texttt{draft\_answers} and \texttt{draft\_saved\_at} via \texttt{gorm.Expr("NULL")}, and posts a "submission received" notification.

\paragraph{Autosave.} \texttt{exam.AutosaveAttempt} (\texttt{internal/exam/handler.go:650--685}) stores the raw JSON body the SPA sends. It refuses to write into a submitted attempt, validates that the body parses as JSON (so a corrupted draft cannot crash the parse later), and stamps \texttt{draft\_saved\_at}.

\section{Auto-grading}

The auto-grader lives in \texttt{internal/judge0}. \texttt{client.go} wraps the Judge0 HTTP API; \texttt{grader.go} contains the three grading functions and the finalisation arithmetic.

\paragraph{Sandbox call.} \texttt{judge0.RunCode} (\texttt{internal/judge0/client.go:64--120}) maps a language string to a Judge0 \texttt{language\_id}, posts the source code and stdin to \texttt{/submissions?wait=true}, and parses the response into a small struct. Two fields are unpacked carefully: the time field is a string in seconds (e.g.\ \texttt{"0.012"}) and is converted to integer milliseconds; the memory field is a JSON \texttt{Number} and is converted to integer kilobytes. Both go onto the \texttt{TestResult} row so the student review screen can show real per-test timings.

\paragraph{Output normalisation.} \texttt{normalizeOutput} (\texttt{internal/judge0/grader.go:366--374}) strips carriage returns, trims trailing whitespace per line, and trims edges. This is the single most important reason auto-grading does not falsely fail students: Judge0 frequently returns \texttt{\textbackslash r\textbackslash n}-terminated output and a stray trailing newline, and a teacher's expected output is rarely typed with bit-perfect line endings.

\begin{figure}[H]
\begin{lstlisting}[language=Go,caption={Output normalisation (\texttt{internal/judge0/grader.go:366--374}).},label={lst:normalize}]
func normalizeOutput(s string) string {
    s = strings.ReplaceAll(s, "\r\n", "\n")
    s = strings.ReplaceAll(s, "\r", "\n")
    lines := strings.Split(s, "\n")
    for i, ln := range lines {
        lines[i] = strings.TrimRight(ln, " \t")
    }
    return strings.TrimSpace(strings.Join(lines, "\n"))
}
\end{lstlisting}
\end{figure}

\paragraph{Coding grader.} \texttt{judge0.Grade} (\texttt{internal/judge0/grader.go:13--109}) loops the test cases for a problem, calls \texttt{RunCode} per case, classifies the result, writes a \texttt{TestResult} row, accumulates max time and memory, and finally writes a single update to the \texttt{Submission} row with the aggregate score and status. \Cref{fig:F5.5} draws the flow.

\begin{figure}[H]
  \centering
  \includegraphics[width=\linewidth]{F5.5_grader-flowchart.pdf}
  \caption{Coding-grader flow (\texttt{internal/judge0/grader.go:13--109}).}
  \label{fig:F5.5}
\end{figure}

The Judge0 \texttt{StatusID} mapping is worth a sentence: \texttt{3} is "Accepted" and \texttt{4} is "Wrong Answer". The grader treats both as a successful run --- the program executed; we just need to compare its output against the expected output ourselves --- and only treats \texttt{5} (TLE), \texttt{6} (compilation error), and \texttt{11} (runtime error) as terminal failures. The reason is that Judge0's idea of "Accepted" depends on the absence of a custom expected output, which APEX always supplies; collapsing 3 and 4 into a single "ran" branch keeps the comparison logic in one place.

\paragraph{MCQ grader.} \texttt{judge0.GradeMCQ} (\texttt{internal/judge0/grader.go:111--210}) parses the \texttt{Problem.CorrectOptionIDs} JSON array and the \texttt{Submission.SelectedOptions} JSON array, builds two sets, and accepts only when the sets are equal. Two corner cases are explicit: an empty correct-set marks the submission \texttt{pending\_review} so a student cannot fail an unconfigured question; an empty selected-set is normalised to \texttt{[]} earlier, in \texttt{exam.SubmitAttempt}, so a "skipped" MCQ never compares as "correct against an empty correct-set". \Cref{fig:F5.6} draws the decision tree.

\begin{figure}[H]
  \centering
  \includegraphics[width=0.9\linewidth]{F5.6_mcq-decision.pdf}
  \caption{MCQ grading decision tree (\texttt{internal/judge0/grader.go:111--210}).}
  \label{fig:F5.6}
\end{figure}

\paragraph{Written grader.} \texttt{judge0.GradeWritten} (\texttt{internal/judge0/grader.go:212--234}) is mostly bookkeeping. It marks the submission \texttt{pending\_review} with \texttt{score = 0} and \texttt{total\_count = 1}, then triggers \texttt{maybeFinalizeAttempt}. Real grading happens later, when the teacher visits the manual-grading queue.

\paragraph{Finalisation.} \texttt{judge0.maybeFinalizeAttempt} (\texttt{internal/judge0/grader.go:245--328}) is the function that turns a set of submissions into one attempt score. It runs at the end of every grader (auto and manual), counts pending/running submissions for the attempt, returns early if any are still in flight, and otherwise iterates the exam's full problem list. Skipped questions stay in the denominator at score zero; \texttt{pending\_review} items are excluded from both numerator and denominator, so a student's aggregate score does not pull down before manual grading lands. \Cref{fig:F5.7} draws the algorithm.

\begin{figure}[H]
  \centering
  \includegraphics[width=\linewidth]{F5.7_finalize-attempt.pdf}
  \caption{Finalisation algorithm (\texttt{internal/judge0/grader.go:245--328}). The \texttt{graded\_notified} flag makes the closing notification idempotent.}
  \label{fig:F5.7}
\end{figure}

\section{Manual grading}

\texttt{internal/submission/handler.go} owns three endpoints: \texttt{RunSolution} (sample-test runs that never persist), \texttt{GradeSubmission} (manual grading by the teacher), and \texttt{GetSubmission} (read-back with per-test detail).

\texttt{submission.GradeSubmission} (\texttt{internal/submission/handler.go:76--128}) is short but security-relevant. It loads the submission, then loads the exam to verify ownership: \texttt{exam.TeacherID == teacherID}. This re-check is the patch documented in \Cref{ch:security}: an earlier revision of the handler trusted the URL alone, so a teacher could grade a submission on a colleague's exam by guessing its ID. The handler then writes the score, calls \texttt{judge0.FinalizeAttempt} so the attempt's aggregate updates, and posts a notification only for standalone (non-exam) submissions --- attempt-bound submissions get the single "Exam Graded" notification from \texttt{maybeFinalizeAttempt}.

\texttt{submission.GetSubmission} (\texttt{internal/submission/handler.go:130--190}) similarly enforces ownership: a teacher may only read submissions on exams they own, and a student may only read their own submissions. The view payload joins per-test detail and selectively reveals the input/expected output for sample tests only --- hidden test cases stay hidden in the response.

\section{Reminder scheduler}

\texttt{reminder.Start} (\texttt{internal/reminder/scheduler.go:16--27}) launches a goroutine that sleeps 15 seconds at boot, runs the sweep once, and then ticks every minute. The sweep itself, \texttt{runOnce} (\texttt{internal/reminder/scheduler.go:29--53}), is a top-level \texttt{recover} wrapper around three calls: \texttt{notifyExam} for in-app reminders within the next 60 minutes, \texttt{emailSweep1Hour} for "starts in about 1 hour" emails, and \texttt{emailSweepStart} for "starting now" emails. Each sweep dedupes via a separate timestamp column, so adding or removing one channel does not affect the others.

The 1-hour and start sweeps use bracketed windows (\texttt{[+58, +62]} minutes and \texttt{[-2, +2]} minutes respectively) so a missed tick is recovered on the next minute. Once a sweep fires, the dedup column is stamped, and subsequent ticks skip the exam. There is no across-process leader election: APEX assumes a single backend instance, and a multi-instance deployment would need to coordinate (a Postgres advisory lock or a Redis-backed lock would suffice).

\section{Frontend SPA}

The SPA is a Vite-bundled React + TypeScript application. \texttt{frontend/src/main.tsx} mounts \texttt{App}; \texttt{App.tsx} wraps a \texttt{BrowserRouter} around \texttt{AuthProvider} and \texttt{AppRoutes}. \Cref{fig:F5.10} draws the provider tree, and \Cref{fig:F5.11} draws the route tree.

\begin{figure}[H]
  \centering
  \includegraphics[width=0.7\linewidth]{F5.10_provider-tree.pdf}
  \caption{Provider tree.}
  \label{fig:F5.10}
\end{figure}

\begin{figure}[H]
  \centering
  \includegraphics[width=\linewidth]{F5.11_route-tree.pdf}
  \caption{Route tree from \texttt{frontend/src/app/routes.tsx}.}
  \label{fig:F5.11}
\end{figure}

\paragraph{AuthContext.} \texttt{frontend/src/contexts/AuthContext.tsx:1--187} owns the session. On mount, it reads \texttt{kernel-token} from \texttt{localStorage}, decodes the JWT in the browser (\texttt{parseJWT}, lines 27--39), and rebuilds the \texttt{AuthUser} object if the token has not expired. A login or signup writes the new token, decodes it, and updates state. The decoded role is mirrored to a separate \texttt{kernel-role} key for backward compatibility with the dashboard layout's guard. Logout drops both keys plus a couple of legacy display-name keys.

The decision to decode the JWT in the browser was deliberate: the token already contains \texttt{user\_id}, \texttt{email}, \texttt{name}, \texttt{role}, and \texttt{exp}, so the session restore avoids one round-trip on every page load. The frontend never \emph{trusts} the decoded payload for authorisation --- every protected operation goes through the API, and the API verifies the signature server-side --- so a client-side tampered token only fools the SPA into showing a wrong dashboard, not into bypassing real checks.

\paragraph{API client.} \texttt{frontend/src/lib/api.ts:1--131} is a thin \texttt{fetch} wrapper that injects the bearer token, handles 401 by clearing the token and redirecting to \texttt{/auth}, and throws a typed \texttt{ApiError} on non-2xx responses. Every endpoint in the rest of the file is a one-line wrapper that posts a typed body and parses a typed response. The convention keeps the rest of the SPA free of \texttt{fetch} calls.

\paragraph{ExamTaking.} The student's exam screen, \texttt{frontend/src/features/exams/pages/ExamTaking.tsx}, is the most complicated component in the SPA. It mounts a Monaco editor for coding answers, a textarea for written answers, and an option list for MCQs. It restores from \texttt{localStorage} or from the server's \texttt{draft\_answers} (whichever is newer), runs a 15-second autosave timer, runs a once-per-second timer to count down to the submit deadline, and posts the final answers on submit. \Cref{fig:F5.12} draws the per-question state.

\begin{figure}[H]
  \centering
  \includegraphics[width=0.85\linewidth]{F5.12_exam-take-state.pdf}
  \caption{ExamTaking state machine.}
  \label{fig:F5.12}
\end{figure}

The autosave snapshot includes the \texttt{answers} array, \texttt{currentIdx}, and \texttt{visitedQuestions} so a fresh device can rehydrate the same view. The component also honours \texttt{exam.reset\_at}: if the exam was destructively edited after the student started, the local snapshot is discarded and the student is forced to re-enter the start screen. This is what makes the exam reset operation safe to use mid-exam.

\paragraph{Other notable pages.} \texttt{ExamBuilder.tsx} is the teacher-side editor: it loads an existing draft via \texttt{getExam}, lets the teacher add coding/MCQ/written problems, and posts them via \texttt{addProblem} and \texttt{addTestCase}. \texttt{GradeWritten.tsx} loads the pending-review queue and lets a teacher set per-question score plus textual feedback. \texttt{ExamReview.tsx} renders a student's attempt with per-test pass/fail tabs and the teacher's feedback, gated by the exam's \texttt{show\_results\_after} setting and the class's \texttt{grades\_announced} flag.

\section{Postgres conventions}

A handful of conventions show up across the SQL surface and are worth flagging in one place.

\paragraph{Idempotent migrations.} Every block in \texttt{RunMigrations} is wrapped in \texttt{IF NOT EXISTS} (for \texttt{ALTER TABLE ADD COLUMN}, \texttt{CREATE INDEX}) or guarded by an \texttt{IS NULL} check (for \texttt{UPDATE} backfills). Replaying the migration step on an already-migrated database is a no-op.

\paragraph{No raw user input in SQL.} Every read goes through GORM's parameterised query helpers (\texttt{Where}, \texttt{First}, \texttt{Find}, \texttt{Pluck}). The only \texttt{db.Exec} calls in the codebase are in the migration step, where they are static SQL strings.

\paragraph{Sentinel cleanups.} A handful of \texttt{UPDATE} statements normalise legacy data: \texttt{submissions.selected\_options} is rewritten from JSON literal \texttt{null} to \texttt{[]}; \texttt{exam.start\_time} is backfilled to \texttt{created\_at} for legacy rows; corrupted \texttt{is\_draft = true} on rows that have a \texttt{start\_time} are flipped to \texttt{false}. Each is a one-line statement gated on the relevant column existing.

\paragraph{Foreign keys after the fact.} The \texttt{problems.folder\_id} foreign key is added in \texttt{RunPostMigrations}, after \texttt{AutoMigrate} has had a chance to create the \texttt{folders} table. The pre-step also sweeps orphaned references to \texttt{NULL} so the constraint can be added without a violation.

\section{Build and run}

The repository ships a one-shot start script, \texttt{start.sh}, that creates \texttt{.env} from \texttt{.env.example} if missing, brings up the Docker Compose Postgres on port 5432, runs \texttt{go run main.go} in the background, waits for \texttt{:8080} to respond, and then runs \texttt{npm install} + \texttt{npm run dev} for the frontend. \texttt{docker-compose.yml} declares the Postgres service with a named volume, so a stop/restart cycle preserves data.

\texttt{frontend/vite.config.ts} declares a dev-server proxy from \texttt{/api} to \texttt{http://localhost:8080}, which is what makes \texttt{frontend/src/lib/api.ts:1} default \texttt{API\_BASE} to \texttt{/api} in development. In production the SPA is built (\texttt{npm run build}) and served as static assets behind the same reverse proxy that fronts the API.

\section{Summary}

The implementation is shaped by the design of \Cref{ch:design} and the constraints of \Cref{ch:requirements}, but two concrete patterns recur across the codebase. Authorisation is checked at three layers --- the JWT middleware (identity), the role middleware (coarse permission), and the handler body (per-resource ownership) --- and each layer fails closed. Idempotency is built into every operation that fires more than once: idempotent migrations, idempotent attempt creation, idempotent reminder dedup, idempotent finalisation. The next chapter looks at where those patterns came from: the security review that pushed the codebase from a working prototype into something defensible.
